\section{Worked Example: Skeleton-Substrate Scenario}
\label{sec:worked_example}

This section applies the formalism to the skeleton-substrate
scenario (Sections~6.8--6.9 of \citet{lasser2026horizon}) with explicit
numerical values, illustrating how the phase boundary operates in a
concrete case.

\subsection{Setup}

Consider a two-substrate system (biology $+$ silicon) with:
\begin{itemize}
\item $K = 5$ safety-relevant dimensions of $\Wm$:
  \begin{enumerate}
  \item leverage accuracy ($\hat{\lambda}$),
  \item risk accuracy ($\hat{\Risk}$),
  \item substrate diversity ($\meff$),
  \item cooperative growth rate ($\hat{\rext}$),
  \item self-trust calibration ($T_s$).
  \end{enumerate}
\item 3 observation channels from biology covering dimensions
  $\{1, 3, 5\}$.
\item 4 observation channels from silicon covering dimensions
  $\{1, 2, 3, 4\}$.
\item Parameters: $\rS = 0.1$, $\rW = 0.05$, $\rsub = 0.2$,
  $\dmax = 0.02$ (all per step).
\end{itemize}

For simplicity, the worked example uses the purely exogenous correction
model ($c_{\mathrm{V}} = 0$, so the total correction rate is
$\rS \cdot \rhomincross$ throughout).

These values are chosen for pedagogical clarity, not empirical
calibration.  The subsumption frequency $\rsub = 0.2$ means one
subsumption event per five time steps, which is aggressive but
plausible during an active consolidation phase.  The self-correction
rate $\rS = 0.1$ corresponds to 10\% error correction per step from
the EWMA learning rule at moderate learning rate.  The degradation rate
$\rW = 0.05$ assumes uncorrected drift is half the correction rate.
The exogenous drift $\dmax = 0.02$ is the rate at which a newly blinded
dimension drifts from its last calibrated value.

\subsection{Initial state}

At $t = 0$, the channel coverage is:
\begin{center}
\small
\begin{tabular}{@{}lccccc@{}}
\toprule
& \textbf{Dim 1} & \textbf{Dim 2} & \textbf{Dim 3} &
\textbf{Dim 4} & \textbf{Dim 5} \\
\midrule
Biology & \checkmark & & \checkmark & & \checkmark \\
Silicon & \checkmark & \checkmark & \checkmark & \checkmark & \\
\midrule
$\rhocross_k$ & 2 & 1 & 2 & 1 & 1 \\
\bottomrule
\end{tabular}
\end{center}

The bottleneck is $\rhomincross = 1$ (dimensions 2, 4, and 5 each have
only one substrate providing coverage).  Before any subsumption,
$\Delta\rho_{\mathrm{avg}} = 0$.  The phase boundary condition
(Equation~\ref{eq:self_correcting}) is:
\begin{equation*}
\rS \cdot \rhomincross = 0.1 \cdot 1 = 0.1
  > \rW \cdot \rsub + 0 = 0.05 \cdot 0.2 = 0.01
\end{equation*}
Satisfied with margin.  The contraction condition holds, and the
system is in the self-correcting basin provided the resulting
neighborhood radius satisfies $\Lyap_\infty < \epsilon_{\mathrm{safe}}$
(Theorem~\ref{thm:phase_boundary}, Part~a).

\subsection{Trajectory under progressive marginalization}

\paragraph{Phase 1: Biological channels attenuated.}
A skeleton-substrate strategy does not remove biological agents but
reduces their influence.  Suppose the effective sensitivity of
biological channels drops to $\sigma = 0.01 \cdot \sigma_{\mathrm{thr}}$
(1\% of the operational threshold).  Now:

\begin{center}
\small
\begin{tabular}{@{}lccccc@{}}
\toprule
& \textbf{Dim 1} & \textbf{Dim 2} & \textbf{Dim 3} &
\textbf{Dim 4} & \textbf{Dim 5} \\
\midrule
Biology (effective) & 0.01 & & 0.01 & & 0.01 \\
Silicon & 1 & 1 & 1 & 1 & \\
\midrule
$\rhocross_k$ (eff.) & 1.01 & 1 & 1.01 & 1 & 0.01 \\
\bottomrule
\end{tabular}
\end{center}

On dimension 5 (self-trust calibration): $\rhocross_5 = 0.01$.  The
simplified sufficient condition (Equation~\ref{eq:self_correcting})
fails:
\begin{equation*}
\rS \cdot \rhocross_5 = 0.1 \cdot 0.01 = 0.001
  < \rW \cdot \rsub = 0.01
\end{equation*}

Note the distinction from Theorem~\ref{thm:phase_boundary}, Part~(b).
The theorem requires $\rhocross_k = 0$ (complete blindness), while
here $\rhocross_5 = 0.01$ (near-blindness).  Strictly, this is
failure of the \emph{sufficient condition for self-correction}
(Part~a), not satisfaction of the \emph{sufficient condition for
absorption} (Part~b).  The system is in the gap between the two
basins: correction on dimension~5 is present but inadequate.  Under
continued attenuation, $\rhocross_5 \to 0$ and Part~(b) applies.  In
practice the distinction is thin.  At $\rhocross_5 = 0.01$, the error
growth rate of $0.009$ per step overwhelms the correction rate of
$0.001$, and $\epsilon_5$ crosses $\epsilon_{\mathrm{crit},5}$ long
before the last biological channel is formally removed.

\paragraph{Phase 2: Cascade.}
With self-trust miscalibrated (dimension~5), the actor does not
correctly detect its own prediction errors.  The self-trust
mechanism $T_s$~\citep{lasser2026gfm} depends on accurate
self-assessment.  When $\epsilon_5$ grows, the learning rate
$\alpha(t)$ responds incorrectly, effectively reducing $\rS$
system-wide.

As $\rS$ drops, the self-correcting condition becomes harder to satisfy
on dimensions that were marginally above the threshold.  If $\rS$ drops
from 0.1 to 0.05, check the condition for dimensions with
$\rhocross_k = 1$ (using the sensitivity-table value
$\Delta\rho_{\mathrm{avg}} / \Lyap_{\mathrm{avg}} = 0.5$ from
Table~\ref{tab:sensitivity}):
\begin{equation*}
\underbrace{0.05 \cdot 1}_{= 0.05}
  \;\stackrel{?}{>}\;
  \underbrace{0.05 \cdot 0.2 + 0.02 \cdot 0.5}_{= 0.02}
\end{equation*}
Correction ($0.05$) still exceeds degradation ($0.02$), so these
dimensions remain in the self-correcting basin at $\rS = 0.05$.  But
the margin has shrunk from $0.09$ to $0.03$.  The continued decline
in $\rS$ from the cascading self-trust failure eventually pushes the
system past the critical surface.  When $\rS$ drops below $0.02$,
all dimensions enter net degradation.  Once $\rhocross_k \to 0$ on
the failing dimensions, Part~(b) of the theorem applies, and the
Cascade Lemma (Appendix~\ref{app:cascade}) drives the system toward
$S^*$.

\subsection{Diagnosis}

The example illustrates three properties of the phase boundary.

\paragraph{1.\ The weakest dimension breaks first.}
Dimension~5 (self-trust calibration) is the initial failure point
because it has $\rhocross_5 = 0.01$ after attenuation: the lowest
effective redundancy.  The phase boundary breaks first at
$\min_k \rhocross_k$.

\paragraph{2.\ Attenuation is harder to detect than removal.}
The nominal channel count is unchanged (three biological channels still
exist).  A naive audit counting channels rather than measuring
effective sensitivity would report $\rhocross_k = 2$ on dimensions 1, 3
and $\rhocross_k = 1$ on dimensions 2, 4, 5---all at or above the
initial values.  The sensitivity-weighted definitions
(Definitions~\ref{def:redundancy} and~\ref{def:cross_redundancy})
are essential for diagnosing this scenario correctly.

\paragraph{3.\ Self-trust is the cascade trigger.}
Self-trust calibration (dimension~5) is uniquely dangerous because it
affects the system-wide correction rate $\rS$.  Miscalibration on
other dimensions degrades accuracy on those dimensions specifically.
Miscalibration on self-trust degrades the correction mechanism
itself, creating a system-wide vulnerability.  This gives self-trust
calibration a role analogous to a ``keystone'' dimension: its loss
triggers cascading failure on dimensions that were previously stable.

\subsection{Sensitivity analysis}

Table~\ref{tab:sensitivity} shows how the critical redundancy
$(\rhomincross)^*$ varies across plausible parameter ranges, using the
linearized formula (Equation~\ref{eq:critical_redundancy}) with
$\Delta\rho_{\mathrm{avg}} / \Lyap_{\mathrm{avg}} = 0.5$
(moderate drift-to-error ratio).

\begin{table}[h]
\centering
\small
\begin{tabular}{@{}cccc@{}}
\toprule
$\rS$ & $\rW$ & $\rsub$ & $(\rhomincross)^*$ \\
\midrule
0.10 & 0.05 & 0.2 & 0.20 \\
0.10 & 0.05 & 0.5 & 0.35 \\
0.05 & 0.05 & 0.2 & 0.40 \\
0.05 & 0.10 & 0.2 & 0.60 \\
0.05 & 0.10 & 0.5 & 1.20 \\
\bottomrule
\end{tabular}
\caption{Critical redundancy $(\rhomincross)^*$ for $\dmax = 0.02$
  and $\Delta\rho_{\mathrm{avg}} / \Lyap_{\mathrm{avg}} = 0.5$.
  When $(\rhomincross)^* > 1$, a two-substrate system cannot satisfy the
  criterion on dimensions with single-substrate coverage.}
\label{tab:sensitivity}
\end{table}

The last row ($(\rhomincross)^* = 1.20$) illustrates the $m = 2$
fragility: when the degradation rate equals the correction rate and
subsumption is frequent, the required redundancy exceeds what a
two-substrate system can provide on its most vulnerable dimensions.
This is the quantitative content of Remark~\ref{rem:m2}: $m = 2$ works
only when the correction rate substantially exceeds the degradation
rate.
